\subsection{The parent objective: increase expected advantage}
\label{subsec:grpo-gradient-direction}

Section~\ref{subsec:grpo-group-advantage} assigns each sampled response an advantage. Hold
those advantages fixed during the subsequent optimizer step and temporarily omit clipping
and reference-policy KL. The conceptual objective is
\begin{equation}
  J_{\mathrm{adv}}(\theta)
  =\mathbb{E}_{o\sim\pi_\theta(\cdot\mid q)}
    \left[A_{\mathrm{old}}(q,o)\right].
  \label{eq:grpo-parent-expected-advantage}
\end{equation}
This is the parent objective from which the reward-driven part of GRPO is constructed. It
is not an extra loss term. It asks a direct question: after changing the policy, will the
policy place more probability on responses that had positive advantage and less on
responses that had negative advantage?

For a two-response example, suppose the behavior policy assigns probability $0.5$ to a
response with advantage $+1$ and probability $0.5$ to a response with advantage $-1$.
Its expected advantage is zero. A new policy assigning probabilities $0.7$ and $0.3$ has
expected advantage
\begin{equation}
  0.7(+1)+0.3(-1)=0.4.
  \label{eq:grpo-parent-objective-example}
\end{equation}
The objective prefers the new policy because probability moved toward the better response.
The advantage therefore supplies the direction of preference; the policy probabilities
determine whether the model actually moves in that direction.

Equation~\ref{eq:grpo-parent-expected-advantage} is easy to state but cannot be evaluated
directly from the current rollout buffer: its expectation is under $\pi_\theta$, while the
stored responses were sampled from $\pi_{\mathrm{old}}$.
Section~\ref{subsec:grpo-importance-ratio} first corrects this sampling-distribution
mismatch exactly with a trajectory importance ratio. It then explains the additional
approximation that produces the per-token $\rho_{i,t}A_i$ term in the GRPO formula.
